\documentclass[UTF8]{ctexart}
\usepackage{amsmath,amssymb,bm}
\title{非线性最小二乘方法}
\author{Serein 回归测试}
\date{2026-07-21}

\begin{document}
\maketitle

\section{模型与残差}

给定观测数据 $\{(\bm x_i,y_i)\}_{i=1}^{n}$，考虑非线性模型

\begin{equation}
y_i=f(\bm x_i;\bm\theta)+\varepsilon_i,
\label{eq:nls-model}
\end{equation}

其中，$\bm\theta\in\mathbb{R}^{p}$ 为待估计参数。残差向量定义为

\begin{equation}
\bm r(\bm\theta)
=
\begin{bmatrix}
r_1(\bm\theta)\\
r_2(\bm\theta)\\
\vdots\\
r_n(\bm\theta)
\end{bmatrix}.
\label{eq:nls-residual}
\end{equation}

\section{目标函数}

非线性最小二乘估计为

\begin{equation}
\hat{\bm\theta}
=
\underset{\bm\theta}{\operatorname{arg\,min}}\,
F(\bm\theta),
\qquad
F(\bm\theta)
=
\frac{1}{2}\bm r^{\mathsf T}(\bm\theta)\bm r(\bm\theta).
\label{eq:nls-objective}
\end{equation}

\subsection{梯度与 Hessian}

\begin{align}
\nabla F(\bm\theta)
&=
\bm J^{\mathsf T}(\bm\theta)\bm r(\bm\theta),
\label{eq:nls-gradient}\\
\nabla^2F(\bm\theta)
&=
\bm J^{\mathsf T}\bm J
+
\sum_{i=1}^{n}r_i(\bm\theta)\nabla^2r_i(\bm\theta).
\label{eq:nls-hessian}
\end{align}

\section{迭代方法}

\subsection{Gauss--Newton 方法}

\begin{equation}
\left(\bm J_k^{\mathsf T}\bm J_k\right)
\Delta\bm\theta_k
=
-\bm J_k^{\mathsf T}\bm r_k.
\label{eq:nls-gn}
\end{equation}

\subsection{Levenberg--Marquardt 方法}

\begin{subequations}
\label{eq:nls-lm-system}
\begin{align}
\left(\bm J_k^{\mathsf T}\bm J_k+\lambda_k\bm I\right)
\Delta\bm\theta_k
&=
-\bm J_k^{\mathsf T}\bm r_k,
\label{eq:nls-lm-step}\\
\bm\theta_{k+1}
&=
\bm\theta_k+\Delta\bm\theta_k.
\label{eq:nls-update}
\end{align}
\end{subequations}

由式~\eqref{eq:nls-lm-step} 可以计算阻尼增量。

\section{终止条件}

\begin{equation}
\left\|\Delta\bm\theta_k\right\|_2
\leq
\varepsilon_{\theta}
\left(\left\|\bm\theta_k\right\|_2+1\right),
\qquad
\left\|\nabla F(\bm\theta_k)\right\|_{\infty}
\leq
\varepsilon_g.
\label{eq:nls-stopping}
\end{equation}

\end{document}
